\section{Probability mass functions}

Once a discrete random variable \(X\) has been defined, we want to know how
probability is distributed among its possible values. This information is encoded
by the probability mass function.

The basic quantity is
\[
P(X=x).
\]
This notation is convenient, but it must be read carefully. The expression
\(X=x\) describes an event in the original sample space:
\[
\{X=x\}=\{\omega\in\Omega:X(\omega)=x\}.
\]
Thus \(P(X=x)\) means the probability of all elementary outcomes that are mapped
to the value \(x\).

\begin{definitionbox}
Let \(X\) be a discrete random variable. Its \textbf{probability mass function}
\(p_X\) is defined by
\[
p_X(x)=P(X=x).
\]
For values outside the support, \(p_X(x)=0\).
\end{definitionbox}

The word mass is used because, in a discrete distribution, probability is
concentrated on separate values. Each value in the support carries a certain
amount of probability mass.

\subsection*{The two basic properties}

A probability mass function must satisfy two conditions:
\[
p_X(x)\geq 0
\]
for every \(x\), and
\[
\sum_{x\in\operatorname{supp}(X)}p_X(x)=1.
\]
These are not arbitrary requirements. They come from the axioms of probability.
The events \(\{X=x\}\) for different support values are disjoint, and their union
is the whole sample space up to events of probability zero.

\begin{remarkbox}
A proposed formula for a probability mass function is not valid until its values
are non-negative and sum to \(1\) over the support.
\end{remarkbox}

\subsection*{Example: indicator of an event}

Let \(A\) be an event and let \(I_A\) be its indicator:
\[
I_A(\omega)=
\begin{cases}
1, & \omega\in A,\\
0, & \omega\notin A.
\end{cases}
\]
The support of \(I_A\) is contained in \(\{0,1\}\). Its mass function is
\[
p_{I_A}(1)=P(I_A=1)=P(A),
\]
and
\[
p_{I_A}(0)=P(I_A=0)=P(A^c)=1-P(A).
\]

Thus an event can be represented as a zero-one random variable. This will be
important when we study expectations and counts.

\subsection*{Example: sum of two dice}

Let \(S(i,j)=i+j\) be the sum of two fair dice. The sample space has \(36\)
equally likely ordered pairs. To find \(p_S(s)\), we count how many pairs have
sum \(s\).

The counts are
\[
\begin{array}{c|ccccccccccc}
s&2&3&4&5&6&7&8&9&10&11&12\\
\hline
\#\{(i,j):i+j=s\}&1&2&3&4&5&6&5&4&3&2&1
\end{array}
\]
Therefore
\[
\begin{array}{c|ccccccccccc}
s&2&3&4&5&6&7&8&9&10&11&12\\
\hline
p_S(s)&\frac1{36}&\frac2{36}&\frac3{36}&\frac4{36}&\frac5{36}&\frac6{36}&\frac5{36}&\frac4{36}&\frac3{36}&\frac2{36}&\frac1{36}
\end{array}
\]

For example,
\[
p_S(7)=P(S=7)=\frac6{36}=\frac16.
\]
The value \(7\) has the largest probability because it can be produced in the
largest number of ordered ways.

\subsection*{Example: number of heads in \(n\) tosses}

Toss a coin \(n\) times independently, with
\[
P(H)=p,\qquad P(T)=1-p.
\]
Let
\[
X=\text{the number of heads}.
\]
The support is
\[
\{0,1,2,\ldots,n\}.
\]
For \(X=k\), exactly \(k\) positions among the \(n\) positions must contain heads.
There are
\[
\binom{n}{k}
\]
ways to choose those positions. Each sequence with \(k\) heads and \(n-k\) tails
has probability
\[
p^k(1-p)^{n-k}.
\]
Hence
\[
p_X(k)=P(X=k)=\binom{n}{k}p^k(1-p)^{n-k},\qquad k=0,1,\ldots,n.
\]

This formula is not just a formula to remember. It contains the whole mechanism:
choose positions of successes, multiply probabilities along each sequence, and
sum over all sequences with the same count.

\subsection*{Example: urn model without replacement}

Suppose an urn contains \(K\) objects of type success and \(N-K\) objects of type
failure. We select \(n\) objects without replacement and ignore order. Let
\[
X=\text{the number of successes selected}.
\]
If \(X=k\), then we choose \(k\) successes from the \(K\) success objects and
\(n-k\) failures from the \(N-K\) failure objects. The total number of possible
samples is \(\binom{N}{n}\). Therefore
\[
p_X(k)=P(X=k)=\frac{\binom{K}{k}\binom{N-K}{n-k}}{\binom{N}{n}},
\]
for values of \(k\) for which the binomial coefficients are meaningful.

This distribution comes from the mechanism of sampling without replacement. It
is different from the binomial model because the selections are not independent:
a success removed from the population changes what remains.

\subsection*{Reading a mass function as a model}

A mass function is a compact description of a probability law. But it should not
be separated from the mechanism that produced it. When reading or constructing a
mass function, ask:
\begin{itemize}
    \item what is the original experiment?
    \item what is the random variable?
    \item what support values are possible?
    \item which elementary outcomes produce each value?
    \item how are their probabilities added?
\end{itemize}

\begin{theorembox}
For a discrete random variable, the probability mass function is obtained by
transferring probability from elementary outcomes to numerical values:
\[
p_X(x)=P(\{\omega\in\Omega:X(\omega)=x\}).
\]
\end{theorembox}
